\chapter{Conclusiones}
\label{cap:conclusiones}

Esta investigación abordó el problema del \textit{bias} y la \textit{fairness} en modelos fundacionales aplicados a \textit{Whole Slide Images} (WSI) de patología digital. El trabajo se sustentó en la premisa de que el buen desempeño agregado de estos modelos no basta para establecer su utilidad clínica: es necesario examinar la composición, procedencia y documentación de los datos, así como su comportamiento ante subgrupos y cambios de distribución. A partir del marco teórico y del estado del arte, se estableció que la variabilidad de las WSI, asociada a la población, los protocolos de preparación, la tinción, la digitalización y la institución de origen, puede afectar la transferencia de las representaciones aprendidas y ocultar disparidades que no son visibles mediante métricas globales \parencite{Basak2023,NEURIPS2024_c9826b9e}.

En respuesta a esta problemática, se revisaron y sistematizaron nueve recursos públicos de patología digital y se construyó un marco operativo para identificar indicadores observables de tres mecanismos potenciales de \textit{bias}: representación, adquisición, y documentación y selección. Posteriormente, se priorizaron seis cohortes de los dominios de cáncer de mama y colorrectal, para las cuales fue posible armonizar variables clínicas, demográficas y técnicas suficientes para definir espacios de evaluación comparables. Sobre esta base, se seleccionaron los modelos Phikon-v2, UNI2-h, Prov-GigaPath y Virchow2 como extractores de representaciones congelados, y se definieron cinco tareas de clasificación orientadas a examinar utilidad predictiva, desempeño por subgrupo y generalización entre cohortes. De este modo, la principal contribución metodológica del estudio consiste en vincular explícitamente la auditoría de los datos con un protocolo de evaluación homogéneo para modelos fundacionales de patología digital.

\paragraph{Objetivo general} El objetivo general, orientado a evaluar distintos niveles de \textit{bias} en modelos fundacionales clínicos mediante una caracterización sistemática de conjuntos de datos de WSI en términos de \textit{fairness} y \textit{bias}, se considera cumplido de manera parcial y en etapa de consolidación. Se completó la construcción de la base conceptual, la caracterización de las cohortes, la selección justificada de modelos y el diseño del protocolo experimental que articula desempeño global, brechas por subgrupo y generalización entre conjuntos de datos. Esta aproximación es coherente con la literatura que advierte que las métricas agregadas pueden ocultar diferencias relevantes entre grupos y que la evaluación de modelos fundacionales debe considerar simultáneamente utilidad y equidad \parencite{NEURIPS2024_c9826b9e}. La evaluación empírica desarrollada en el Capítulo~6 confirma esta aproximación: el riesgo de \textit{bias} documentado por dataset en la Tabla~\ref{tab:matriz_riesgo_unificada} predijo, de forma más consistente que la elección de modelo fundacional, tanto el nivel de utilidad alcanzado como la estabilidad de las conclusiones de equidad en las cinco tareas evaluadas (Sección~\ref{sec:sintesis_cap6}). Por lo tanto, el objetivo general se considera cumplido: se logró evaluar sistemáticamente distintos niveles de \textit{bias} en modelos fundacionales clínicos, articulando la caracterización de datos con un protocolo experimental común.

\paragraph{Objetivo específico 1} El primer objetivo específico, consistente en caracterizar conjuntos de datos de imágenes médicas de patología digital mediante la identificación de variables demográficas y parámetros de adquisición críticos para el estudio de \textit{bias}, se considera cumplido. La revisión de nueve recursos públicos permitió evidenciar una marcada heterogeneidad en tamaño de cohorte, estructura y completitud de metadatos. Tras el proceso de curación y armonización, se identificó que la edad constituye el atributo demográfico más viable para el análisis estratificado en los dominios seleccionados, mientras que el sexo no es utilizable como atributo sensible en cáncer de mama debido a la homogeneidad casi absoluta de las cohortes. Asimismo, se observó una asociación estructural entre la cohorte de origen y los parámetros técnicos de adquisición, tales como escáner, magnificación y tinción, lo cual limita la atribución de diferencias de desempeño exclusivamente a las variables clínicas. Este hallazgo es consistente con la evidencia que señala que la procedencia institucional puede dejar firmas visuales en imágenes histopatológicas y favorecer el aprendizaje de señales ajenas a la morfología de interés \parencite{alsaafin2023biased,howard2021site}. La caracterización desarrollada no demuestra por sí sola una brecha de \textit{fairness}, pero establece las condiciones necesarias para interpretar con cautela los resultados experimentales posteriores.

\paragraph{Objetivo específico 2} El segundo objetivo específico, referido a seleccionar modelos fundacionales representativos y definir tareas específicas acordes con los conjuntos de datos analizados, se considera cumplido. La selección de Phikon-v2, UNI2-h, Prov-GigaPath y Virchow2 se fundamentó en criterios de pertinencia para patología computacional, compatibilidad con WSI teñidas con hematoxilina y eosina, accesibilidad de pesos e implementación, documentación técnica, trazabilidad del preentrenamiento, diversidad metodológica y viabilidad computacional. Al emplear los cuatro modelos como extractores de representaciones congelados y mantener constante el procesamiento posterior, el protocolo reduce la influencia de adaptaciones específicas y favorece una comparación controlada de las representaciones preentrenadas. Además, la definición de tareas de clasificación de estado HER2, subtipo molecular, sitio anatómico, lateralidad y órgano permite examinar escenarios clínicos y anatómicos con distintos niveles de comparabilidad entre cohortes. Esta decisión se alinea con la literatura sobre modelos fundacionales en patología, que destaca la transferibilidad de representaciones auto-supervisadas, pero también la necesidad de evaluar sus supuestos de preentrenamiento y trazabilidad \parencite{chen2024uni,filiot2024phikonv2,xu2024provgigapath,vorontsov2024virchow}.

\paragraph{Objetivo específico 3} El tercer objetivo específico, correspondiente a evaluar el desempeño de los modelos seleccionados mediante métricas acordes con las tareas definidas y analizar el impacto del \textit{bias} sobre su capacidad de generalización, se encuentra en desarrollo. El protocolo ya establece una evaluación a nivel de paciente que incluye métricas de utilidad global, análisis por clase, brechas entre subgrupos, experimentos \textit{cross-dataset} e intervalos de confianza obtenidos mediante \textit{bootstrap} no paramétrico. Este diseño permite distinguir entre desempeño promedio, estabilidad de las estimaciones y variación asociada al cambio de cohorte, evitando interpretar una única métrica agregada como evidencia suficiente de robustez o equidad. De esta forma, se evaluaron cuatro modelos fundacionales en cinco tareas de clasificación, reportando utilidad global, brechas por subgrupo etario y de sexo, y generalización \textit{cross-dataset} con intervalos de confianza por \textit{bootstrap}. Las caídas de desempeño ante cambio de distribución fueron sustanciales y transversales (16--73\% de pérdida relativa en T1--T4), y el experimento dedicado de verificación del \textit{center effect} (Sección~\ref{sec:verificacion_center_effect}) cuantificó que una parte sustancial de estas caídas es atribuible a la firma de adquisición e institución, y no exclusivamente a un cambio morfológico real. El objetivo específico 3 se considera, por tanto, cumplido.

\paragraph{Hipótesis de investigación} La hipótesis planteó que los conjuntos de datos de imágenes médicas presentan \textit{bias} demográficos y de adquisición significativos, los cuales afectan negativamente la capacidad de generalización de modelos fundacionales para WSI y producen diferencias de desempeño entre subgrupos poblacionales. La evidencia disponible permite respaldar parcialmente su primer componente: la caracterización identificó desbalances de representación, restricciones de cobertura de metadatos y heterogeneidad técnica sistemáticamente asociada a las cohortes. Estos antecedentes son compatibles con mecanismos de \textit{bias} de representación y adquisición, y con el riesgo de \textit{domain shift} y \textit{shortcut learning} descrito en la literatura \parencite{suresh2021framework,mehrabi2021survey,geirhos2020shortcut,stacke2021domainshift}. Así, los sesgos de representación y adquisición documentados a lo largo del proyecto afectan la generalización, tal como confirma el experimento de verificación del \textit{center effect} (BA$\approx$0.85 entre cohortes de mismo escáner, BA=1.000 al introducir un escáner distinto), pero no respalda el segundo componente en las cohortes de tamaño adecuado: las brechas de equidad por edad y sexo fueron bajas y sin un grupo sistemáticamente desfavorecido en BCNB, HISTAI-Breast, SurGen y HISTAI-CRC-B1 ($<$0.10 en la mayoría de los casos). Las brechas elevadas observadas en HSI-BC e HISTAI-CRC-B2 no permiten distinguirse de ruido muestral dado el tamaño de esas cohortes y la amplitud de sus intervalos de confianza. En consecuencia, la hipótesis se considera aceptada parcialmente: se confirma que existen sesgos de adquisición con efecto medible sobre la generalización, pero no se encuentra evidencia de disparidades de desempeño sistemáticas entre subgrupos demográficos dentro del rango de datos evaluado.

\paragraph{Trabajo futuro}

El principal aprendizaje metodológico de este trabajo es que la evaluación de fairness y generalización en modelos fundacionales no depende únicamente de la elección del encoder, de la tarea predictiva o de las métricas empleadas. Depende de las características y limitaciones de los conjuntos de datos disponibles. La curación y armonización realizadas permitieron construir escenarios comparables para cáncer de mama y colorrectal; sin embargo, también evidenciaron que gran parte de los conjuntos de datos públicos de patología digital presenta documentación heterogénea, cobertura incompleta de variables clínicas y demográficas, tamaños muestrales desiguales y condiciones de adquisición estrechamente asociadas a cada cohorte.

Por otra parte, las diferencias de generalización entre cohortes no pueden atribuirse de forma aislada a un único factor. Cada dataset reúne simultáneamente variaciones en institución de origen, escáner, preparación de la muestra, protocolo de tinción, población clínica, criterios de anotación y proceso de curación de metadatos. La alta separabilidad del origen de las cohortes observada en los embeddings confirma que los modelos retienen información asociada al dataset; sin embargo, no permite determinar qué proporción de esta señal corresponde a características técnicas, clínicas o morfológicas reales.

A partir de ello, una línea de trabajo futuro prioritaria consiste en ampliar la disponibilidad y calidad de los datos antes de proponer mecanismos de mitigación a nivel de modelo. Sería especialmente valioso contar con cohortes multicéntricas que reporten de manera sistemática variables demográficas, atributos clínicos, detalles de adquisición, procedimientos de tinción, escáner utilizado y criterios de anotación. Además, sería necesario incorporar muestras que permitan desacoplar parcialmente estas fuentes de variación; por ejemplo, cohortes provenientes de distintas instituciones que compartan protocolos técnicos, o muestras obtenidas con distintos escáneres dentro de una misma institución.

También sería relevante evaluar los modelos sobre cohortes de mayor tamaño y con una representación más equilibrada de subgrupos demográficos y clases clínicas. Esto permitiría estimar métricas de fairness con menor incertidumbre, analizar posibles intersecciones entre atributos, por ejemplo, edad y sexo, y estudiar si las brechas se mantienen al modificar la composición de la población evaluada. En este contexto, la falta de una disparidad consistente en las cohortes analizadas no debe interpretarse como demostración de equidad clínica, sino como evidencia de que las diferencias observables están condicionadas por el soporte y la calidad de los datos disponibles.

Finalmente, se podrían evaluar estrategias de adaptación o mitigación, tales como normalización de tinción, adaptación de dominio, fine-tuning supervisado o aprendizaje federado \citep{Xu2025StainNormalization,Kumari2024DomainAdaptation,Zhang2024FederatedLearning}. No obstante, estas estrategias deberían evaluarse sobre datos suficientemente documentados y bajo escenarios externos bien definidos.